\scanpage{39}
\margin{\unsure{NB.} S'il s'agit de \unsure{grps} de type galoisien, et $U$ est ouvert dans $I$, alors il s'ensuit qu'il contient un ss.-grp ouvert \unsure{ds} $U_{\unsure{0}}$, qui est \uline{invariant} dans $D$.}

\uline{\unsure{Sort|Suite} de la \unsure{moyenne}}

$D$ grp., $I$ \unsure{ss.-grp} distingué, $E$ vectoriel de dim finie sur corps $k$, sur lequel $D$ opère. On suppose que $I$ admet un sous-groupe d'indice fini $U$, \inserted{invariant dans $D$,} tel que $g\in U\Rightarrow g$ opère sur $E$ de façon \uline{unipotente}, i.e. valeurs propres \unsure{tous|toutes} $1$. Alors le groupe algébrique d'autom. de $E$ engendré par les $g\in U$ est ``unipotent'', donc \deleted{\illegible{2}} $E$ admet une suite de composition sur $U$ telle que les \inserted{\unsure{$E$ opère} \unsure{triviale|trivialem.} sur les} facteurs \deleted{soient \illegible{1} sur $E$}. De façon précise, on peut prendre la suite de composition par les $E^U$, $(E/E^U)^U$ $\ldots$. Cette suite est stable par les opérations de $D$, \uline{si} \deleted{(comme il est loisible)} \deleted{on prend} \inserted{puisque} $U$ dans $I$ invariant par tout $D$. Donc il existe une suite de composition de $E$ \uline{sur $D$} tels que sur chaque facteur $E_i$, \deleted{$I$} $U$

\margin{38}

